\documentclass[10pt]{article}

%%% These are some packages that are useful
\usepackage{amsmath,amssymb, amscd,amsbsy, amsthm, enumerate}
\usepackage[export]{adjustbox}
\usepackage{lastpage}
\usepackage[top=1in, bottom=1in, left=1in, right=1in]{geometry}
\usepackage[unicode]{hyperref}
\usepackage{tikz, pgfplots, xcolor, fancyhdr}
\usepackage{multicol,caption}
\usepackage{lipsum}
\usepackage[version=4]{mhchem}
\usepackage{float}
\usepackage{enumitem}
\usepackage{xurl}

%%% Page formatting
%\setlength{\headsep}{30pt}
\setlength{\textheight}{9in}
\newcommand{\tab}{\hspace{1cm}}
%\setlength{\parindent}{25pt}

\title{Know Thy Enemy (Part 6) 	: A Guard Against Possession}
\author{Antonius Torode}

%%% Header and Footer Info
\pagestyle{fancy}
\fancyhead[L]{{\large Template - \textbf{Change 03}}}
\fancyhead[C]{\today}
\fancyhead[R]{Name: Antonius Torode}


\fancyhf{} % sets both header and footer to nothing
\renewcommand{\headrulewidth}{0pt}
% your new footer definitions here

\fancyfoot[L]{}
\fancyfoot[C]{}
\fancyfoot[R]{\thepage\ of \pageref{LastPage}}

% Used to define spacing and format of References
\let\OLDthebibliography\thebibliography
\renewcommand\thebibliography[1]{
	\OLDthebibliography{#1}
	\setlength{\parskip}{0pt}
	\setlength{\itemsep}{0pt plus 0.3ex}
}

\newenvironment{Figure}
  {\par\medskip\noindent\minipage{\linewidth}}
  {\endminipage\par\medskip}


%%% Document Starts now
\begin{document}

\maketitle
\thispagestyle{fancy}


%\begin{abstract}
%TODO
%\end{abstract}

\begin{multicols}{2}

``Build your house upon the rock'' is a well known phrase among Christians. It is a concept that can be understood from many different perspectives and is a cornerstone of some of the most important Biblical lessons.
\begin{quotation}
``Therefore whoever hears these sayings of Mine, and does them, I will liken\footnote{Strong's 3666 [homoioó]: Meaning to compare or resemble \cite{Bible Hub}.} him to a wise man who built his house on the rock\footnote{Strong's 4073 [petra]: Meaning a rock, large stone, or cliff \cite{Bible Hub}.}: and the rain descended, the floods came, and the winds blew and beat on that house; and it did not fall, for it was founded on the rock.'' - Matthew 7:24-25
\end{quotation}
Have you ever considered what this concept means from the perspective of demonic possession? That question may sound wildly out of place, however these two things can actually be seen as two sides of the same coin. Possession \textit{could} be described as the rain descending, the floods coming, or the winds blowing and beating down on someone's house. If that's the case, could building your house on a rock actually be the key to \textit{protecting yourself} from demonic possession?

Taking a step back, what does scripture mean by your \textit{house}? In the context of possessions, it literally means \textit{you}.
\begin{quotation}
``Do you not know that you are the \textit{temple}\footnote{Strong's 3485 [naos]: Means a shrine, temple or sanctuary (a divine dwelling-place) \cite{Bible Hub}.} of God and that the Spirit of God \textit{dwells}\footnote{Strong's 3611 [oikeó]: Means to occupy a house or to figuratively reside \cite{Bible Hub}.} in \textit{you}?'' - 1 Corinthians 3:16
\end{quotation}
The scriptures continually use the metaphor that \textit{you} are akin to a house (a spiritual dwelling), and if you live God's way, he will reside \textit{with you}.
\begin{quotation}
``If anyone loves Me, he will keep My word; and My Father will love him, and We will come to him and make Our \textit{home} with him.'' - John 14:23
\end{quotation}

Next, what is the rock, and what does it mean to build yourself upon that rock? 
\begin{quotation}
``and all drank the same spiritual drink. For they drank of that spiritual \textit{Rock} that followed them, and that Rock was \textit{Christ}.'' - 1 Corinthians 10:4
\end{quotation}
The \textit{Rock} is representative of Christ. Building yourself upon him means to use his teachings as a foundation for your life. It means to build up His character within yourself - to take in and accept his spiritual drink. This in turns opens yourself (your house), so that he can figuratively reside inside you. If you don't do this, the rain will descend and the floods will come. What exactly does that mean?

Scriptures speak of a concept called possession, which has many examples. As English speakers, we have a particular understanding of this.
\begin{quotation}
Definition [possession]: ``controlled or \textit{overwhelmingly influenced} by something (such as an evil spirit, a passion, or an idea) \cite{possession_defn}.''
\end{quotation}
In scriptures, there are two relevant ways (that I'll mention here) that possession is referred to. The first is two separate Greek words, d\textit{aimonion echó}, which translates to \textit{have a demon}. The word \textit{echó} here means to have, to hold, or to possess - such as being in possession of. That sounds almost like the demon is there with you, in your house.
\begin{quotation}
``For John came neither eating nor drinking, and they say, `He \textit{has a demon}\footnote{Strong's 2192 + 1140 [daimonion echó]: Meaning to have, possess, or hold an evil spirit (demon).}.'\thinspace'' - Matthew 11:18
\end{quotation}
Another common term is \textit{daimonizomai}, which more literally means \textit{to be demon possessed} or \textit{to be under the power of a demon}. The definition of this Greek term is interesting though, because it also means to \textit{have a devil(s)} or to be \textit{vexed with a devil(s)}. Now, having a demon and being vexed is an interesting way to look at possession, because it does not really imply that the demon has taken control of you. It instead implies that there is a demon with you, and he's essentially just harassing you. The definitions of vexed \cite{vexed_defn} follow as:
\begin{itemize}
	\setlength\itemsep{0em}
	\item `feeling or showing irritation, annoyance, or distress'
	\item `difficult and often frustrating to deal with'
	\item  `affected with, marked by, or causing trouble or harassment.'
\end{itemize}
If we combine all of these concepts, we can see that being demon possessed is essentially the same as a demon residing in \textit{you} (or your \textit{house}). From there, he's meddling, irritating, annoying, distressing, and much more. Even the old testament describes evil spirits in this way:
\begin{quotation}
``But the Spirit of the Lord departed from Saul, and a distressing spirit from the Lord \textit{troubled}\footnote{Strong's 1204 [baath]: meaning to terrify, to startle, to dismay.} him. And Saul’s servants said to him, “Surely, a distressing spirit from God is \textit{troubling} you.'' - 1 Samual 16:14-15
\end{quotation}
Imagine being harassed by the same entity day and night because they are literally there with you in your home at all times. This is why the English definition of possession says \textit{overwhelmingly influenced by}. That is possession.

In order to protect ourselves from it, we must build ourselves on the Rock. Christ is that Rock we must build our foundation on. That is because He is the one with power over the evil spirits and the power to cast them out. 
\begin{quotation}
``When evening had come, they brought to Him many who were \textit{demon-possessed}\footnote{Strong's 1139 [daimonizomai]: Meaning to be demon-possessed, to be under the power of a demon, to have a demon(s), or to be vexed with demon(s) \cite{Bible Hub}.}. And He \textit{cast out} the spirits with a word'' - Matthew 8:16
\end{quotation}
When we build our house on that Rock, he sends his spirit to dwell in us and those evil spirits flee from us\footnote{James 4:7 - ``Therefore submit to God. Resist the devil and he will flee from you.''}. So then what happens if we do not build out house upon the Rock? Can Christ not just cast out those evil spirits as he did to the others. Would we then be safe from the floodwaters and evil spirits?
\begin{quotation}
``When an unclean spirit goes out of a man, he goes through dry places, seeking rest, and finds none. Then he says, ‘I will return to my \textit{house} from which I came.’ And when he comes, he finds it \textit{empty}, swept, and put in order.  Then he goes and takes with him seven other spirits more wicked than himself, and they enter and \textit{dwell} there; and the last state of that man is worse than the first.'' - Matthew 12:43-45
\end{quotation}
Just as Christ offers to dwell in your home, so too can the evil spirits. If we leave our house empty, we leave an open invitation to any spirit that wants to come make itself at home. 

To take it one step further, this concept does not just apply to you as an individual. It applies to any home. It applies to churches, to cities, and even to entire nations. As nations move away from God, their foundation turns from Rock to sand. Then the rains descend, and the winds blow, and they beat upon the house. That is the evil spirits moving in and beginning to dwell. The nation becomes possessed. The only protection against this is to build your house upon the Rock.


\begin{quotation}
``But everyone who hears these sayings of Mine, and does not do them, will be like a foolish man who built his house on the sand:  and the rain descended, the floods came, and the winds blew and beat on that house; and it fell. And great was its fall.'' - Matthew 7:
26-27
\end{quotation}
That fall is merely the evil spirits moving in to their new home.


\begin{thebibliography}{9}
{\footnotesize
	\bibitem{Bible Hub} \url{https://biblehub.com}
	
	\bibitem{possession_defn} ``Possessed,'' Merriam-Webster.com Dictionary, Merriam-Webster, \url{https://www.merriam-webster.com/dictionary/possessed}. Accessed Oct 17 2026.
	
	\bibitem{vexed_defn} ``Vexed.'' Merriam-Webster.com Dictionary, Merriam-Webster, \url{https://www.merriam-webster.com/dictionary/vexed}. Accessed Oct 29 2026.
}
\end{thebibliography}

\end{multicols}


%%%%%%%%%%%%%%%%%%%%%%%%%%%%%%%%%%%%%%%%%%%%%%%%%%%%%%%%%%%%%%%%%%%%%%%%%%%%%%%%%%%%%%%%%%%
\end{document}





















}{den}